% ──────────────────────────────────────────────────────────────────────────────
% Capitolo 1 – Introduzione
% ──────────────────────────────────────────────────────────────────────────────
\chapter{Introduzione}
\label{ch:introduzione}

\section{Conformità ai Vincoli di Progetto e Riproducibilità}
\label{sec:vincoli_progetto}
Il presente sistema, denominato SINTON-IA, è stato architettato e strutturato in rigorosa rispondenza ai vincoli accademici previsti:
\begin{enumerate}
    \item \textbf{Problema Reale}: SINTON-IA affronta le criticità operative insite nelle applicazioni di \textit{e-mental health}: l'incapacità umana di filtrare manualmente e proattivamente enormi flussi testuali quotidiani alla ricerca di intenti suicidari latenti, l'incongruenza delle risposte nei questionari di self-report psicologico e l'altissimo tasso di abbandono terapeutico preventivo (\textit{churn rate}) tipico di queste piattaforme prima della reale presa in carico.
    \item \textbf{Sperimentazione Multi-Pipeline}: La ricerca comprende il training valutato di almeno due algoritmi paralleli per ogni funzionalità esplorata (es. TF-IDF + Logistic Regression vs Reti Neurali per NLP; LightGBM vs Random Forest per regressione; Algoritmo Genetico vs Random Search per l'ottimizzazione del Churn).
    \item \textbf{Simulazioni ed Estensione Dati ad-hoc}: È stato interamente sviluppato in Python un simulatore probabilistico ad-hoc che genera iterazioni comportamentali per l'addestramento dell'Algoritmo Genetico, mentre i dataset linguistici da fonti esterne hanno subito profondi rimaneggiamenti strutturati di pre-processing e test sintetico di robustezza contestuale. 
    \item \textbf{Metodologia e Trade-Off Computazionali}: L'intera fase di Model Training (Capitolo \ref{ch:modeling}) è corredata da metriche ad hoc (F1-score pesati $F_2$, MAE bilanciati, curve di convergenza ad Early Stopping) con discussioni formali ed esplicite sui Trade-off scelti dal team (Performance predictive bilanciate contro Tempi di Esecuzione sub-secondo e complessità dei calcoli hardware).
    \item \textbf{Output Fruibile}: La validazione operativa del sistema viene garantita dall'integrazione dei moduli di intelligenza artificiale sviluppati direttamente all'interno dell'infrastruttura software preesistente di SINTONIA, dimostrandone l'immediata applicabilità clinica.
\end{enumerate}

Per garantire il libero accesso alla documentazione tecnica, ai notebook esplorativi, ai file Python (\texttt{.py}) e alla gestione dell'ambiente, l'intero ciclo di vita applicativo è integralmente tracciato tramite controlli di versione. Il codice sorgente prodotto risulta costantemente accessibile e consultabile presso il \textbf{Repository GitHub Ufficiale}: \\ \url{https://github.com/AntonioWalter/SINTON-IA}

\clearpage
\section{Il Problema dell'Accesso alle Cure Psicologiche in Campania}
\label{sec:contesto_campania}
Nel contesto attuale della sanità pubblica campana, l'accesso ai servizi di supporto psicologico risulta spesso complesso, caratterizzato da lunghe liste d'attesa che possono protrarsi per diversi mesi prima dell'effettiva presa in carico da parte di uno psicologo dell'ASL. Durante l'attesa, i cittadini non dispongono di strumenti o servizi digitali che possano sostenerli. Questa mancanza unita al forte distacco istituzionale provoca spesso un crescente senso di abbandono e una percezione di scarsa attenzione da parte del servizio sanitario pubblico.

Negli ultimi anni, in particolar modo a causa delle conseguenze post-pandemiche, la domanda di assistenza psicologica è aumentata in modo significativo, mentre la capacità di risposta delle strutture pubbliche, limitata dalla carenza di personale e fondi, è rimasta per molti versi invariata. Questa situazione evidenzia la profonda assenza di strumenti automatizzati che possano assistere i cittadini durante i mesi di attesa e migliorare la gestione delle richieste in carico alle strutture.

In risposta a queste criticità è nata \textbf{SINTONIA}, una piattaforma applicativa innovativa \cite{sintonia_repo}, concepita specificatamente per ridurre l'impatto psicologico scaturito dall'attesa e migliorare l'efficienza dei processi di presa in carico del cittadino, promuovendo un approccio di cura preventiva e inclusiva. SINTONIA rappresenta un passo concreto verso un sistema di supporto psicologico pubblico più accessibile, moderno, e tempestivo sul territorio Campano.

\clearpage
\section{La Piattaforma SINTONIA e il Monitoraggio Attivo}
\label{sec:sintonia_app}
Per comprendere il valore e il contesto di fondo del progetto \textit{SINTON-IA}, risulta cruciale comprendere le funzionalità e gli obiettivi dell'ecosistema sorgente.

SINTONIA risolve il problema delle estenuanti liste di attesa integrando un software di monitoraggio continuo per la salute mentale. Tramite l'impiego di questionari clinici e di autovalutazione certificati \cite{kroenke2001phq9} (es. PHQ-9, GAD-7, PC-PTSD-5), l'applicazione è in grado di ricalcolare costantemente l'urgenza di ogni paziente, generando una priorità clinica e una lista di attesa dinamica in grado di favorire l'intervento tempestivo a supporto dei casi più gravi.

Oltre a una periodica raccolta di questionari, SINTONIA introduce ulteriori avanzati sistemi di monitoraggio attivo e passivo:
\begin{itemize}
    \item \textbf{Diario personale}: uno spazio libero dove i pazienti possono trascrivere pensieri e sensazioni quotidiane;
    \item \textbf{Stato d'animo giornaliero}: la capacità di quantificare, con valenza e intensità, l'andamento del proprio umore;
    \item \textbf{Forum clinico e comunitario}: piattaforme comunicative sicure in cui i pazienti possono condividere riflessioni ed estrarre supporto o ricevere rapide risposte da psicologi certificati.
\end{itemize}

L'abbondanza dei dati comportamentali e testuali raccolti tramite queste metriche periodiche istituisce un perfetto \textit{ground truth} e pone solide basi in vista dell'integrazione di sistemi decisionali e predittivi basati sull'Intelligenza Artificiale (SINTON-IA). Questi saranno predisposti a semplificare ulteriormente il lavoro degli psicologi (analizzando moli di dati irrealizzabili per un occhio umano) e a supportare autonomamente il cittadino in attesa della terapia.

\clearpage
\section{Il Flusso di Lavoro dello Psicologo: Criticità e Limiti Attuali}
\label{sec:flusso_psicologo}

Nella versione attuale del sistema \cite{sintonia_repo}, lo psicologo dispone di una vista centralizzata attraverso cui amministra i pazienti presi in carico. Sebbene la piattaforma garantisca una gestione strutturata, l'interazione pratica e il monitoraggio quotidiano presentano limitazioni operative significative:

\vspace{0.4cm}
\begin{itemize}
    \item \textbf{Revisione manuale dispendiosa:} Lo psicologo è costretto ad effettuare controlli e \textit{revisioni a campione casuale} dei questionari per scovare false compilazioni, distrazioni o palesi incongruenze. Si tratta di una procedura farraginosa che comporta un pesante confronto incrociato manuale con lo storico clinico, rallentando drasticamente il processo di \textit{triage}.

    \item \textbf{Barriere di privacy nel monitoraggio:} A tutela della privacy pre-presa in carico, il professionista non ha l'autorizzazione per leggere il contenuto integrale del \textit{diario personale} del paziente. L'osservazione testuale diretta rimane rilegata unicamente alle porzioni pubbliche del sistema, come il forum anonimo.

    \item \textbf{Gestione degli Alert disconnessa dal tempo reale:} Il sistema genera allarmi (Alert Clinici) quando le risposte a un questionario (es. PHQ-9) delineano un profilo di rischio gravissimo. L'enorme criticità risiede tuttavia nella \textit{frequenza} di questi test, somministrati a intervalli regolari (es. ogni 14 giorni). Questo schema temporale rende il sistema "congelato" nei restanti giorni: se il paziente esplicita intenti autolesivi nel proprio diario nei giorni che intercorrono tra un test e un altro, l'allarme non ha alcun modo di innescarsi.
\end{itemize}

Questi evidenti vuoti di granularità temporale e di carico cognitivo sullo psicologo gettano le basi per l'integrazione di \textit{SINTON-IA}. Un livello di algoritmi intelligenti capace di analizzare in background e in \textit{real-time} i dati non supervisionati, fornendo allarmi proattivi indipendenti dalle scadenze dei questionari convenzionali.

% ──────────────────────────────────────────────────────────────────────────────
% Sezione 1.2 – Funzionalità Implementate
% ──────────────────────────────────────────────────────────────────────────────
\clearpage
\section{Integrazione AI: Funzionalità Previste}
\label{sec:funzionalita_implementate}
Il progetto SINTON-IA prevede l'integrazione di tre componenti di Intelligenza Artificiale, ciascuno progettato per colmare una specifica lacuna operativa dell'attuale piattaforma SINTONIA. In questa fase introduttiva ci si limita a descrivere la finalità e il tipo di task affrontato da ciascun componente, rimandando ai capitoli successivi l'analisi tecnica e architetturale.

\subsection{Red Flag Detection (Rischio Suicidario)}
\label{subsec:red_flag}
Il primo modello affronta un task di \textbf{classificazione su testo} ed è concepito per analizzare in modo completamente anonimo le pagine scritte dai pazienti all'interno del diario personale. L'obiettivo principale è l'individuazione automatica di pattern linguistici riconducibili a un potenziale rischio suicidario.

Qualora il modello rilevi la presenza di tali pattern, il paziente verrà classificato come soggetto a rischio e il sistema genererà automaticamente un \textit{Alert Clinico}, notificando tempestivamente lo psicologo di riferimento. Questa funzionalità risulta particolarmente strategica poiché opera su una fonte dati che, per ragioni di privacy, non è direttamente consultabile dai professionisti prima della presa in carico.
In questo modo il modello colma il vuoto temporale degli alert basati esclusivamente sui questionari, estendendo la copertura del monitoraggio a ogni giorno in cui il paziente scrive nel proprio diario.

\subsection{Depression Prediction}
\label{subsec:depression}
Il secondo modello si configura come un task di \textbf{regressione} e si basa sullo storico degli stati d'animo registrati quotidianamente dal paziente all'interno dell'applicazione. L'obiettivo è stimare il punteggio atteso nella compilazione del questionario PHQ-9 (Patient Health Questionnaire), uno degli strumenti clinici più utilizzati per la valutazione della severità della depressione.

Le predizioni generate dal modello saranno impiegate con una duplice finalità:
\begin{itemize}
    \item \textbf{Grafici di andamento}: dopo la presa in carico, lo psicologo potrà visualizzare l'andamento stimato dell'umore del paziente nel tempo, ottenendo una visione continua e granulare del decorso emotivo, utile sia per la diagnosi sia per il monitoraggio terapeutico.
    \item \textbf{Rilevazione di compilazioni incongruenti}: confrontando il punteggio predetto dal modello con quello effettivamente ottenuto nella compilazione del questionario, il sistema potrà evidenziare automaticamente discrepanze significative. Questo meccanismo supporterà concretamente gli psicologi nell'identificazione di false compilazioni, riducendo drasticamente il carico della revisione manuale a campione.
\end{itemize}

\subsection{Churn Prevention}
\label{subsec:churn}
La prevenzione dell'abbandono dell'applicazione (\textit{churn}) da parte dei pazienti durante la fase di attesa viene affrontata mediante un \textbf{approccio di tipo genetico-evolutivo}.

L'obiettivo è stimare la probabilità che un utente smetta di utilizzare attivamente SINTONIA prima della presa in carico, evento che comprometterebbe la raccolta dei dati necessari al corretto funzionamento del triage dinamico. Sulla base delle stime prodotte, il sistema invierà \textit{notifiche strategiche e personalizzate} agli utenti a rischio di abbandono, incentivandoli a mantenere un utilizzo costante della piattaforma.

Il Churn Prevention riveste un ruolo trasversale: garantendo la continuità nell'acquisizione dei dati comportamentali, contribuisce indirettamente all'efficacia degli altri due modelli e al funzionamento complessivo del sistema di prioritizzazione clinica.

